\documentclass{article}

\usepackage[utf8x]{inputenc}
\usepackage[english, russian]{babel}
\usepackage{graphicx}
\usepackage{amsmath}
\usepackage{amssymb}
\usepackage{extarrows}
\usepackage{vmargin}
\usepackage{MnSymbol}
\usepackage{cases}
\setpapersize{A4}
\setmarginsrb{2cm}{2cm}{2cm}{2cm}{0pt}{0mm}{0pt}{13mm}

%user commands
\newtheorem{theorem}{Теорема}
\newtheorem{lemma}{Лемма}

\newenvironment{proof}{\paragraph{Доказательство:}}{\hfill$\blacksquare$}
\newenvironment{definition}{ \paragraph{Определение:}}{\hfill $\blacktriangleleft$}
\newenvironment{observation}{ \paragraph{Замечание:}}{}
\newenvironment{hence}{ \paragraph{Следствие:}}{}

\begin{document}


\centerline{\huge Теормин по курсу} 
\centerline{\huge \textbf{ЧИСЛЕННЫЕ МЕТОДЫ В ИНТЕГРАЛЬНЫХ }}
\centerline{\huge \textbf{УРАВНЕНИЯХ И ИХ ПРИЛОЖЕНИЯ}}
\centerline{\huge для магистров ВМК МГУ}

\centerline{[проверялся, но могут встречаться ошибки]}

\vspace{0.5cm}



\section{Интегралы и интегральные операторы}



\begin{theorem}
(Критерий интегрируемости) $\int_{a}^{b}f(x)dx$ существует $\Leftrightarrow$ 
\begin{enumerate}
\item $f(x)$ - ограничена;
\item множество точек разрыва функции имеет меру Лебега равную 0.
\end{enumerate}
\end{theorem}
\vspace{0.2cm}

\noindent
Пусть $D\in \mathbb{R}^n$, точка $c \in D$.
$F_c(x)$- функция с особенностью в точке $c$, $F_c(x) \geqslant 0$.
Пусть 
$$I = \lim_{d \to 0} \int_{D \backslash U(c)} F_c(x)dx,$$ 
$$J = \sup_{U(c)} \int_{D \backslash U(c)} F_c(x)dx.$$
\begin{theorem}
$I$ и $J$ существуют/не существуют одновременно, и $I = J$, если существуют. 
\end{theorem}

\begin{theorem}
(Мажорантный признак сходимости) Пусть $F_c(x)$ и $G_c(x)$ - функции с особенностью в точке $c$ и $0 \leqslant G_c(x) \leqslant F_c(x)$. Тогда если $\exists \int_D F_c(x)dx = I$, то $\exists \int_D G_c(x)dx = J$, причём $J \leqslant I$.
\end{theorem}

\begin{theorem}
(Об абсолютной интегрируемости) Пусть $F_c(x)$ не обязательно знакопостоянна. Тогда, если $\exists \int_D |F_c(x)|dx$, то $\exists \int_D F_c(x)dx$
\end{theorem}

\noindent
Пусть $D$ - замыкание области (или кривая, или поверхность) в $\mathbb{R}^n$, $D$ - ограничена и замкнута. Пусть $\varphi(x)$ - функция определённая на $D$. 
\begin{definition}
\textbf{Интегральным уравнением Фредгольма I рода} называется:
\begin{equation}
f(x) = \int_D K(x,y) \varphi(y)dy, \quad x \in D.
\end{equation}

\begin{itemize}
\item $K(x,y) \in C[D \times D]$ - с непрерывным ядром;
\item $K(x,y) = \dfrac{K^*(x,y)}{|x-y|^{\alpha}}$ - с интегральным ядром с полярной особенностью, где $K^*(x,y) \in C[D \times D]$, и $\alpha < \dim D$. 
\end{itemize}
\end{definition}

\begin{definition}
Оператор $K(x,y)$ называется \textbf{ограниченным}, если $\exists c: ||f|| \leqslant c ||\varphi||$, $f = K\varphi$.
\end{definition}

\begin{definition}
\textbf{Нормой оператора} $K(x,y)$ называется $||K|| = \sup_{\varphi \neq 0} \dfrac{||K\varphi||}{||\varphi||} = \sup_{||\varphi|| = 1} \dfrac{||K\varphi||}{||\varphi||}$.
\end{definition}

\begin{theorem}
Если $K(x,y)$ - интегральный оператор с непрерывным ядром, то $K: C(D) \to C(D)$ и ограничен.
\end{theorem}


\section{Элементы теории приближений}



\section{Численные методы решения интегральных уравнений Фредгольма 2-го рода}



\section{Численное решение сингулярных интегральных уравнений}



\end{document}